
% Usual aneurysm intro, need for new methods (faster for medical routine)
Intracranial aneurysms (IAs) are pathological dilations of cerebral arteries, most frequently found in the Circle of Willis, that affect approximately 2–3\% of the population. These aneurysms pose a significant clinical challenge due to their potential to rupture, resulting in subarachnoid hemorrhage, a condition with a case fatality rate over 30\% \cite{vlak_prevalence_2011, nieuwkamp_changes_2009}. Current clinical decisions regarding aneurysm management rely heavily on static, macroscopic metrics such as size, shape, and location, along with patient-specific factors like age and hypertension history \cite{baharoglu_identification_2012}. However, these metrics have limited reliability in predicting rupture risk, often leading to inconclusive or contradictory assessments \cite{xiang_hemodynamicmorphologic_2011}.

Recent advances in computational methods and imaging technologies have enabled the modeling of aneurysmal blood flow, highlighting hemodynamic biomarkers such as wall shear stress (WSS), oscillatory shear index (OSI), and intra-aneurysmal pressure gradients as key determinants of rupture risk \cite{xiang_hemodynamicmorphologic_2011}. Computational fluid dynamics (CFD) has been instrumental in providing high-fidelity, patient-specific hemodynamic simulations \cite{zhou_association_2017, xiang_cfd_2014}. Despite its utility, the prohibitive computational cost, time-intensive setup, and specialized expertise required for CFD limit its integration into clinical workflows, where timely decision-making is critical.

In parallel, artificial intelligence (AI) and machine learning (ML) methods have demonstrated transformative potential across a range of neurovascular interventions, as evidenced by their application in device tracking during endovascular procedures. For instance, real-time AI-assisted platforms like Neuro-Vascular Assist (iMed technologies, Tokyo, Japan) have improved precision and reduced cognitive load during interventions such as aneurysm coiling and liquid embolization \cite{sakakura_realtime_2024, sakakura_pioneering_2024}, demonstrating the clinical feasibility of AI-driven decision support. Inspired by these successes, researchers have begun exploring ML-based surrogates for CFD, aiming to achieve accurate, near-instantaneous predictions of hemodynamics \cite{arzani_machine_2022}. Traditional ML approaches have largely relied on convolutional neural networks (CNNs) due to their effectiveness in image-based tasks. While CNNs have shown promise in predicting certain hemodynamic parameters from 2D or simplified 3D data \cite{gharleghi_transient_2022, ferdian_wssnet_2022, rutkowski_enhancement_2021}, their reliance on structured grids makes them ill-suited for the complex, unstructured nature of aneurysm geometries. This limitation is particularly patent when attempting to model the intricate, three-dimensional flow patterns critical to understanding aneurysm pathophysiology. 

Graph neural networks (GNNs) provide a compelling alternative by leveraging the graph-like structure of CFD domains, enabling the direct processing of irregular, non-Euclidean data. Going beyond point cloud-based methods like PointNet and PointNet++\cite{charles_pointnet_2017, qi_pointnet_2017} which have been tested on simplified geometries and stationary hemodynamics predictions\cite{li_prediction_2021_1, suk_physics-informed_2024}, GNNs rely on message passing mechanisms to iteratively update unstructured graphs, and efficiently learn local patterns.  

Recent studies have demonstrated the efficiency of GNNs in predicting 2D time-dependent flows, steady-state 3D flows, and external flows around simple geometries, particularly in physics-driven tasks involving partial differential equations and numerical simulation \cite{pfaff_learning_2021, hernandez_thermodynamics-informed_2024, lam_graphcast_2023, li_physics-constrained_2024, gladstone_gnn-based_2023, pelissier_graph_2024}. To the best of our knowledge, while GNNs have been used in biomedical blood flow and hemodynamic studies \secondround\cite{sen_physics_2024, suk_mesh_2024, pegolotti_learning_2024, feiger_accelerating_2020, li_prediction_2021_2, liang_feasibility_2020}\color{black}, their application to time-dependent 3D blood flow predictions with intricate vortex and swirling patterns within intracranial aneurysms remains an unexplored frontier. 

This work introduces a novel framework that integrates CFD-trained, physics-constrained graph neural networks (GNNs) to predict full 3D, time-resolved hemodynamic fields for intracranial aneurysms throughout the cardiac cycle. To enhance the accuracy of capturing complex spatio-temporal dynamics, the framework incorporates physics-based embedding laws and constraints directly into the learning process. Additionally, it leverages detailed node feature enhancements within the graph representation, enabling precise modeling of intricate hemodynamic variations. Therefore, the proposed learning-prediction framework achieves a balance between computational efficiency and accuracy, addressing the need for high-fidelity simulations while meeting the immediacy required for clinical applications. \firstround To further validate the generalizability of our approach, we perform two inference experiments on out-of-distribution cases: new inflow profiles unseen during training, and four fully patient-specific aneurysm geometries. These results highlight the model’s robustness and its potential as a transferable foundation for future clinical extensions. \color{black} Finally, to support the broader research community, we introduce a comprehensive dataset comprising 105 patient-specific aneurysm geometries coupled with high-resolution CFD simulations, referred to as BenchAnXplore. Unlike existing datasets, which often suffer from limited size or heterogeneity, this dataset enables standardized evaluation of ML models in aneurysm modeling. 

This work lays the foundation for future clinical applications with significant potential impact. Indeed, real-time prediction of hemodynamics could support future advances in rupture risk assessment, enable more personalized treatment planning \cite{hachem_reinforcement_2023}, assist in optimizing interventions such as flow diverters or coiling \cite{chong_multicenter_2019}, and accelerate long-term simulations of intracranial aneurysm healing \cite{pelissier2025continuum}. Moreover, the introduction of a comprehensive benchmark dataset helps address a major limitation in the field, the scarcity of large-scale, high-quality training data for ML models. By integrating modern ML architectures with domain-specific knowledge, this study contributes a foundational step toward a new generation of AI-assisted neurovascular modeling. It lays the groundwork for future innovations in clinical decision support, with the long-term goal of improving outcomes for patients with intracranial aneurysms.
